\scanpage{223}
\mspage{759}
\diagram{\begin{tikzcd}[ampersand replacement=\&,column sep=small,row sep=small]
(\Pi \arrow[r,phantom,"\subset" description] \& \mathcal{G} \arrow[r,phantom,"\subset" description] \& M), \\
\Pi_0 \arrow[u,phantom,"\cup" description] \arrow[r,phantom,"\subset" description] \& \mathcal{G}_0 \arrow[u,phantom,"\cup" description] \& {}
\end{tikzcd}}
and \uline{if} $\mathcal{G}_0\xrightarrow{\varphi}\underline{\mathcal{G}}_{\mathbb Q}^{\deleted{\unsure{\varphi}}}$ comes from a Lie homomorphism of this system, $M\longrightarrow\underline{M}(\unsure{A|\mathbb{A}}_{\mathbb Q})$, by the method of the preceding section, then we find a \unsure{natural} homomorphism
\diagram{\begin{tikzcd}[ampersand replacement=\&,column sep=large]
(71)\quad\underline{M}_{\mathbb Q} \arrow[r] \& \underline{M}_{\mathbb Q}^{\varphi},
\end{tikzcd}}
\margin{via the operation of $\underline{M}_{\mathbb Q}$ on $\underline{\mathcal{G}}_{\mathbb Q}$, and $\chi_{\underline{M}}$}
making the diagram commutative
\diagram{\begin{tikzcd}[ampersand replacement=\&,column sep=large,row sep=large,/tikz/execute at end picture={
\begin{scope}[overlay]
\draw ([xshift=-4mm,yshift=-3mm]\tikzcdmatrixname-1-1.south west) -- ([xshift=4mm,yshift=5mm]\tikzcdmatrixname-1-1.north east);
\draw ([xshift=-5mm,yshift=-1mm]\tikzcdmatrixname-2-1.north west) -- ([xshift=1mm,yshift=2mm]\tikzcdmatrixname-1-2.north west);
\draw ([xshift=-3mm,yshift=-1mm]\tikzcdmatrixname-2-1.south west) -- ([xshift=6mm,yshift=0mm]\tikzcdmatrixname-1-2.south east);
\draw ([xshift=-3mm,yshift=-2mm]\tikzcdmatrixname-3-1.south west) -- ([xshift=7mm,yshift=1mm]\tikzcdmatrixname-2-2.south east);
\draw ([xshift=0mm,yshift=-4mm]\tikzcdmatrixname-3-1.south east) -- ([xshift=5mm,yshift=2mm]\tikzcdmatrixname-3-2.north east);
\end{scope}
}]
\deleted{\mathcal{G}_0} \arrow[r] \arrow[d] \& \deleted{M} \arrow[d] \& {} \& \underline{\mathcal{G}}_{\mathbb Q} \arrow[r] \arrow[dr,"\mathrm{can.}","{\begin{gathered}\text{via \unsure{automorphisms}}\\\mathrm{int.}\end{gathered}}"'] \& \underline{M}_{\mathbb Q} \arrow[d] \\
(72)\quad\deleted{\underline{\mathcal{G}}(\mathbb Q)} \arrow[r,no head] \arrow[d,"{\deleted{\underline{\mathrm{int}}}}"'] \arrow[dr] \& \deleted{\underline{M}(\unsure{A|\mathbb{A}}_{\mathbb Q})} \arrow[d] \& {} \& {} \& \underline{M}_{\mathbb Q}^{\varphi}. \\
\deleted{\underline{M}_{\mathbb Q}^{\varphi}(\mathbb Q)} \arrow[r,hook] \& \deleted{\underline{M}_{\mathbb Q}^{\varphi}(\unsure{A|\mathbb{A}}_{\mathbb Q})} \& {} \& {} \& {}
\end{tikzcd}}

If the center of $\underline{\mathcal{G}}_{\mathbb Q}$ is trivial (which is a ``credible'' hypothesis especially if the center of $\mathcal{G}_0$ is so\ldots) then $\underline{\mathcal{G}}_{\mathbb Q}\longrightarrow\underline{M}_{\mathbb Q}^{\varphi}$ is injective, and $\underline{\mathcal{G}}_{\mathbb Q}$ is interpreted as ``the $\underline{\mathcal{G}}_{\mathbb Q}$'' associated to a Lie homomorphism
\diagram{\begin{tikzcd}[ampersand replacement=\&,column sep=large]
{\left(\begin{array}{ccccc}
\unsure{\pi|\Pi}&\subset&\mathcal{G}&\subset&M\\[-1pt]
\cup&&\cup&&\\[-1pt]
\Pi_0&\subset&\mathcal{G}_0&&
\end{array}\right)} \arrow[r] \& (\underline{\Pi}\subset\underline{\mathcal{G}}_{\mathbb Q}\subset\underline{M}_{\mathbb Q}^{\varphi}),
\end{tikzcd}}
which has a certain universal character with respect to the given homomorphism (65) $\varphi:\mathcal{G}_0\longrightarrow\underline{\mathcal{G}}(\mathbb Q)$.
\margin{222}

\scanpage{224}
When we do not assume in advance the existence of a \deleted{$\unsure{M|H}$} Lie homomorphism in which $\underline{\mathcal{G}}_{\mathbb Q}$ fits, it is not automatic that we can find a homomorphism
\diagram{\begin{tikzcd}[ampersand replacement=\&,column sep=large]
(73)\quad\varphi_M:\ M \arrow[r] \& \underline{M}_{\mathbb Q}^{\varphi}(\unsure{A|\mathbb{A}}_{\mathbb Q})
\end{tikzcd}}
which accounts for the operations of $M$ on $\mathcal{G}_{\deleted{\unsure{0|\mathbb Q}}}$, through the commutativity of a diagram, for $\forall\ g\in M$
\diagram{\begin{tikzcd}[ampersand replacement=\&,column sep=large,row sep=large,/tikz/execute at end picture={
\begin{scope}[overlay]
\draw ([xshift=-1mm,yshift=-1mm]\tikzcdmatrixname-1-1.east) -- ([xshift=-8mm,yshift=-2mm]\tikzcdmatrixname-1-2.west);
\end{scope}
}]
\deleted{M} \arrow[r] \& \mathcal{G} \arrow[r,"{\operatorname{int}(g)|\mathcal{G}}"] \arrow[r,shorten >=12pt] \arrow[d,"\widehat{\varphi}"'] \& \mathcal{G} \arrow[d,"\widehat{\varphi}"] \\
(74) \& \underline{\mathcal{G}}(\unsure{A|\mathbb{A}}_{\mathbb Q}) \arrow[r,"{\operatorname{int}(\varphi_M(g))|\underline{\mathcal{G}}(\unsure{A|\mathbb{A}}_{\mathbb Q})}"'] \& \underline{\mathcal{G}}(\unsure{A|\mathbb{A}}_{\mathbb Q}).
\end{tikzcd}}

We are mainly interested in those $\varphi$ (65) for which (73) exists --- at least after possibly restricting to an open \unsure{subgroup|subgroup} of $M$\ldots
\margin{223}

\scanpage{225}
\mspage{760}

\uline{53}

$1^\circ)$ \unsure{Differentiable|Differential} action and special Teichmüller group.

Let $X$ be a connected \inserted{\deleted{\unsure{compact}} differentiable,} orientable surface (compact in the case of interest to us), $I$ a discrete subset of $X$. For $\forall i\in I$, let $\mathbb{S}_i$ be the space (homeomorphic to the \deleted{\unsure{c|d}} \unsure{circle}) of homogeneous half-lines in the tangent space, and let $\widetilde{\mathbb{S}}_i$ be a universal covering of $\mathbb{S}_i$.

Let \deleted{\unsure{Let} $\unsure{\widetilde g}$} $\unsure{u}$ be a differentiable automorphism of $(X,S)$, it induces an automorphism of
\[
(1)\qquad \mathbb{S}=\mathbb{S}_I\overset{\mathrm{d\acute ef}}{=}\coprod_{i\in I}\mathbb{S}_i,
\]
denoted $u_{\mathbb{S}}$. Let us consider the set of liftings \inserted{$\widetilde u$} of $u_{\mathbb{S}}$ to an automorphism of \deleted{$\unsure{\mathbb{S}}$ of $\unsure{\widetilde{\mathbb{S}}}$}
\[
(2)\qquad \widetilde{\mathbb{S}}=\coprod_{i\in I}\widetilde{\mathbb{S}}_i.
\]
Let us note, if $T_0$ is the orientation module of $X$, that \deleted{each} $T_0$ operates on each $\widetilde{\mathbb{S}}_i$,

\margin{224}

\scanpage{226}

so that $T_0$ in fact operates on $\mathbb{S}$ and
\[
(3)\qquad \widetilde{\mathbb{S}}/T_0\simeq\mathbb{S}.
\]
This shows that (if $u$ is as above) the group $T_0^I$ operates on the left (for example) on the set of liftings $\widetilde u$\unsure{.} which thus becomes a torsor under $T_0^I$.

Let $G$ be a topological group, which operates \unsure{differentiably} on $(X,S)$, with \deleted{\unsure{of|some}} \unsure{one} \inserted{\unsure{of the}} continuous actions on $\mathbb{S}$. We set
\[
(4)\qquad \widetilde G=
\left\{(\deleted{\unsure{g}}\unsure{u},\widetilde u)\;\middle|\;
\begin{array}{l}
u\in G,\quad\widetilde u\text{ a continuous lifting}\\
\text{of }u_{\mathbb{S}}\text{ to }\widetilde{\mathbb{S}}
\end{array}
\right\},
\]
we find on $\widetilde G$ a canonical topology \inserted{(convergence of $u$ and uniform convergence of $\widetilde u$)}, which is such that the homomorphism $G\longrightarrow\widetilde G$ \unsure{makes|will make} $G$ a covering of $\widetilde G$, with group $\unsure{T}_0^I$:
\[
(6)\qquad
\diagram{\begin{tikzcd}[ampersand replacement=\&,column sep=large]
1\arrow[r,no head] \& T_0^I\arrow[r] \& \widetilde G\arrow[r] \& G\arrow[r] \& 1
\end{tikzcd}}.
\]
\margin{\uline{N.B.} $G$ operates on $T_0^I$ via the character \deleted{$\unsure{(52)|(\varepsilon)}$} of orientation}

\uline{Proposition} \deleted{\unsure{If}} Suppose that $X$ is \unsure{compact|compact},

\margin{225}

\scanpage{227}
\mspage{761}

and consider the canonical homomorphism
\[
\diagram{\begin{tikzcd}[ampersand replacement=\&,column sep=large]
G\arrow[r] \& G/G^\circ\arrow[r] \& \mathcal{T}(X,S)
\end{tikzcd}}
\]
(Teichmüller group \inserted{of $(X,S)$}), we then have a canonical homomorphism of extensions
\[
(7)\qquad
\diagram{\begin{tikzcd}[ampersand replacement=\&,column sep=large,row sep=large]
1\arrow[r,no head] \& T_0^I\arrow[r]\arrow[d,"1"] \& \widetilde G\arrow[r]\arrow[d] \& G\arrow[r,no head]\arrow[d] \& 1\\
1\arrow[r] \& T_0^I\arrow[r] \& \mathrm{Sp}\mathcal{T}(X,S)\arrow[r] \& \mathcal{T}(X,S)\arrow[r,no head] \& 1
\end{tikzcd}},
\]
$\mathrm{Sp}\mathcal{T}(X,S)$ is the \uline{special} Teichmüller group of $(X,S)$.

\margin{For the proof cf. the developments of $5^\circ)$ (pages $766'$ ff, \unsure{in particular} pages \unsure{784} ff)}

\uline{N.B.} We shall note that if for $\forall i$ \inserted{$\in I$}, we fix an origin $e_i$ on $\mathbb{S}_i$, then it follows from the universal property of $\widetilde{\mathbb{S}}$ that giving a lifting $\widetilde u$ of $u_{\mathbb{S}}$ to $\widetilde{\mathbb{S}}$ amounts to giving, for $\forall i\in I$, a lifting of $u_{\mathbb{S}}(e_i)$ to a point $\widetilde a_i$ of $\widetilde{\mathbb{S}}_{\unsure{u(e_i)|u(i)}}$ (i.e., of $\widetilde{\mathbb{S}}_{u(i)}$).

\margin{226}

\scanpage{228}

$2^\circ)$ \uline{The groups} $\mathrm{GL}(2,\mathbb{R})^{\pm1}$\unsure{${}^{\sim}$|and} \deleted{$\unsure{\mathrm{GL}(2,\mathbb{Z})}$}. \inserted{We set}
\[
(8)\qquad \mathrm{GL}(2,\mathbb{R})^{\pm1}
=\{u\in\mathrm{GL}(2,\mathbb{R})\mid\det u\in\{\pm1\}\}.
\]
This group operates \unsure{differentiably} on $\mathbb{C}\simeq\mathbb{R}^2$ leaving the origin fixed, we can therefore apply the construction \unsure{of} $1^\circ)$, Here $\mathbb{S}_0$ is identified with the unit circle
\[
(9)\qquad
\diagram{\begin{tikzcd}[ampersand replacement=\&,column sep=large]
\mathbb{U}\overset{\mathrm{d\acute ef}}{=}\{z\in\mathbb{C}\mid|z|=1\}\arrow[r] \& \mathbb{S}_0^{\unsure{\deleted{*}|*}}\\
z\arrow[r,mapsto] \& \mathbb{R}^{+}z
\end{tikzcd}}
\]
and for $u\in\mathrm{GL}(2,\mathbb{R})^{\unsure{\pm1}}$, its action on $\mathbb{S}_0^{\unsure{\deleted{*}|*}}$ is identified with the action on $\mathbb{U}$ given by
\[
(10)\qquad
\diagram{\begin{tikzcd}[ampersand replacement=\&,column sep=large,row sep=0.3em]
z\arrow[r,mapsto] \& \dfrac{u(z)}{|z|}\\
\mathbb{U}\arrow[u,phantom,"\in"{rotate=90}]
\end{tikzcd}}.
\]
A lifting of this action \inserted{$u_{\mathbb{S}_0^{\unsure{\deleted{*}|*}}}$} to an action on $\widetilde{\mathbb{S}}_0^{\unsure{\deleted{*}|*}}$ amounts to a lifting of $u(\unsure{e_0}=1)$ to an element of $\widetilde{\mathbb{S}}_0$. \deleted{\unsure{and}}
We \deleted{\unsure{shall choose|shall describe} the covering} \inserted{\unsure{shall construct|construct} $\widetilde{\mathbb{S}}$ \deleted{$\unsure{\widetilde{\mathbb{S}}_0}$}} by choosing a base point $e_0$ of $\mathbb{S}_0$, corresponding to the

\margin{227}